% Fichier généré automatiquement par tools/generate_corrections.py.
% Source : RDM/RDM-03-Torsion/539_Cours_RDM/539_Cours_RDM.tex
% Statut : partiel (3/5 question(s) renseignée(s), 0 reprise(s)).
\begin{corrigebox}[Corrigé partiel extrait de la banque]
\CorrigeQuestion{1}
$\tau = \dfrac{M_t}{I_G}r$, $\tau = G r \gamma_x$, $M_t= G \gamma_x I_G$ et $\Delta \theta_x = \gamma_x L$.
\CorrigeQuestion{2}
On a $\int_G(S)=\int_S r^2 \dd S$.

Pour un cylindre plein : on a $I_G(S)=\dfrac{\pi D^4}{32} = \dfrac{\pi R^4}{2}$.
Pour une poutre creuse : $I_G(S) = \dfrac{\pi \left(R^4 -r^4\right) }{2} $.
\CorrigeQuestion{3}
\CorrigeACompleter{Aucun élément de réponse n'est présent dans la source d'origine pour cette question.}
\CorrigeQuestion{4}
\CorrigeACompleter{Aucun élément de réponse n'est présent dans la source d'origine pour cette question.}
\CorrigeQuestion{5}
L'angle unitaire de torsion est donné par $\gamma_x = \dfrac{M_t }{GI_G}$ avec $G=\dfrac{E}{2(1+\nu)}$.

AN : 
$E = \SI{200 000}{MPa}$ et $\nu = 1$, alors $G = \SI{77 000}{MPa}$
$I_{G_p} = \dfrac{\pi 5^4}{2} = \SI{981}{mm^4}$ 

$\gamma_x = \dfrac{M_t }{I_G G} = \dfrac{10 000}{77 000\times 981} = \SI{1,3e-4}{rad}$
\end{corrigebox}
